\documentclass[aspectratio=169,10pt]{beamer}
\usetheme[numbering=fraction,progressbar=frametitle]{metropolis}
\usepackage{fontspec}
\setsansfont{Latin Modern Sans}
\usepackage{amsmath,amssymb,mathtools,graphicx,booktabs}
\usepackage{polyglossia}
\setdefaultlanguage{french}
\setbeamertemplate{navigation symbols}{}
\setbeamerfont{frametitle}{size=\large}
\setbeamerfont{normal text}{size=\normalsize}
\setlength{\parskip}{.45em}
\hypersetup{pdfauthor={},pdfsubject={Lecture slides},pdfkeywords={},pdftitle={Calcul intégral}}
\providecommand{\N}{\mathbb{N}}
\providecommand{\Z}{\mathbb{Z}}
\providecommand{\D}{\mathbb{D}}
\providecommand{\Q}{\mathbb{Q}}
\providecommand{\R}{\mathbb{R}}
\providecommand{\C}{\mathbb{C}}
\providecommand{\vect}[1]{\overrightarrow{#1}}
\providecommand{\vi}{\vec{\imath}}
\providecommand{\vj}{\vec{\jmath}}
\providecommand{\vk}{\vec{k}}
\providecommand{\norm}[1]{\left\lVert #1\right\rVert}
\providecommand{\vc}[2]{\begin{pmatrix}#1\\#2\end{pmatrix}}
\providecommand{\deter}[2]{\det\left(#1,#2\right)}
\providecommand{\abs}[1]{\left\lvert #1\right\rvert}
\providecommand{\dd}{\,\mathrm{d}}
\providecommand{\e}{\mathrm{e}}
\providecommand{\Card}{\operatorname{Card}}
\providecommand{\E}{\operatorname{E}}
\providecommand{\V}{\operatorname{V}}
\providecommand{\degre}{^{\circ}}
\providecommand{\un}[1]{\,\mathrm{#1}}
\providecommand{\sci}[2]{#1\times10^{#2}}
\providecommand{\pgcd}{\operatorname{PGCD}}
\providecommand{\ppcm}{\operatorname{PPCM}}
\title{Calcul intégral}
\subtitle{Cours et exemples}
\date{}
\author{}
\begin{document}
\begin{frame}[plain]
\vfill
{\usebeamerfont{title}\inserttitle\par}
\medskip
{\usebeamerfont{subtitle}\insertsubtitle\par}
\medskip
{\usebeamercolor[fg]{progress bar}\rule{\textwidth}{.5pt}}
\vfill
\end{frame}
\begin{frame}{Objectifs}
Calculer et encadrer une intégrale ; distinguer aire et intégrale ; utiliser une intégration par parties ; étudier une suite d'intégrales.
\end{frame}
\begin{frame}{Aire et primitive}
Pour \(f\) continue et positive sur \([a,b]\), \(a\le b\), l'intégrale \(\int_a^b f(x)\,\dd x\) est l'aire sous la courbe, au-dessus de l'axe, entre \(x=a\) et \(x=b\), exprimée en unités d'aire.
\end{frame}
\begin{frame}{Aire et primitive}
Si \(f\) est continue sur un intervalle \(I\) et \(a\in I\), alors \(F_a(x)=\int_a^x f(t)\,\dd t\) est la primitive de \(f\) qui s'annule en \(a\). Pour toute primitive \(F\),
\[\int_a^b f(x)\,\dd x=F(b)-F(a).\]
Toute fonction continue sur un intervalle admet des primitives. Une intégrale est un nombre ; une primitive est une fonction.
\end{frame}
\begin{frame}{Aire et primitive}
Pour \(h>0\), les rectangles donnent \(hf(x)\le F_a(x+h)-F_a(x)\le hf(x+h)\). Après division par \(h\) et passage à la limite, le taux d'accroissement tend vers \(f(x)\). Pour \(h<0\), la croissance donne \(f(x+h)\le f(t)\le f(x)\) sur \([x+h,x]\). En intégrant,
\[(-h)f(x+h)\le F_a(x)-F_a(x+h)\le(-h)f(x).\]
En divisant par \(-h>0\), on obtient
\[f(x+h)\le\dfrac{F_a(x+h)-F_a(x)}h\le f(x).\]
Les deux bornes tendent vers \(f(x)\). Par continuité, \(F_a'=f\). Toute autre primitive est \(F_a+C\), d'où la formule de calcul aux bornes.
\end{frame}
\begin{frame}{Propriétés}
Pour des fonctions continues et des bornes appartenant à un même intervalle :
\[\int_a^a f=0,\quad\int_b^a f=-\int_a^b f,\quad
\int_a^c f=\int_a^b f+\int_b^c f,\quad
\int_a^b(\lambda f+\mu g)=\lambda\int_a^b f+\mu\int_a^b g.\]
Si \(a\le b\) et \(f\le g\), alors \(\int_a^b f\le\int_a^b g\). En particulier, si \(m\le f\le M\), alors \(m(b-a)\le\int_a^b f\le M(b-a)\).
Pour une fonction de signe quelconque, l'intégrale est définie par les primitives. L'aire géométrique entre la courbe et l'axe est \(\int_a^b|f|\), obtenue en découpant aux changements de signe. Entre deux courbes, elle vaut \(\int_a^b|f-g|\).
\end{frame}
\begin{frame}{Propriétés}
\(\int_{-1}^{1}x\,\dd x=0\), mais l'aire entre la droite et l'axe sur cet intervalle vaut \(1\).
\end{frame}
\begin{frame}{Valeur moyenne et intégration par parties}
Pour \(a<b\), la valeur moyenne est \(\mu=\dfrac1{b-a}\int_a^b f(x)\,\dd x\). Le rectangle de hauteur \(\mu\) et de largeur \(b-a\) possède la même aire algébrique.
\end{frame}
\begin{frame}{Valeur moyenne et intégration par parties}
Si \(u,v\) sont de classe \(\mathcal C^1\) sur \([a,b]\),
\[\int_a^b u(x)v'(x)\,\dd x=[u(x)v(x)]_a^b-\int_a^b u'(x)v(x)\,\dd x.\]
\end{frame}
\begin{frame}{Valeur moyenne et intégration par parties}
Intégrer \((uv)'=u'v+uv'\) puis isoler l'intégrale cherchée. Le crochet est indispensable.
\end{frame}
\begin{frame}{Valeur moyenne et intégration par parties}
Avec \(u=x\) et \(v'=e^x\), \(\int_0^1xe^x\,\dd x=[xe^x]_0^1-\int_0^1e^x\,\dd x=e-(e-1)=1\).
\end{frame}
\begin{frame}{Rectangles et suites d'intégrales}
Pour \(f\) continue et croissante sur \([a,b]\), avec \(a<b\) et \(n\ge1\) et \(h=\dfrac{b-a}{n}\), les sommes \(h\sum_{k=0}^{n-1}f(a+kh)\) et \(h\sum_{k=1}^{n}f(a+kh)\) encadrent l'intégrale. Leur différence est \(h(f(b)-f(a))\).
Pour \(I_n=\int_0^1x^nf(x)\,\dd x\), commencer par le signe et les comparaisons de l'intégrande. Si \(0\le f\le M\), alors \(0\le I_n\le \dfrac{M}{n+1}\) ; les gendarmes donnent \(I_n\to0\). Une intégration par parties peut aussi fournir une relation entre \(I_n\) et \(I_{n+1}\).
\end{frame}
\begin{frame}{Rectangles et suites d'intégrales}
\begin{columns}[T,onlytextwidth]
\begin{column}{.57\textwidth}
Dans \(I_n=\int_0^1x^nf(x)\,\dd x\), l'expression \(x^0\) désigne la fonction constante \(1\), y compris à la borne \(0\). Ainsi \(I_0=\int_0^1f(x)\,\dd x\) ; aucun calcul indéterminé n'est nécessaire.   L'aire entre \(y=x\) et \(y=x^2\) sur \([0,1]\) vaut \(\dfrac{1}{6}\).
\end{column}
\begin{column}{.4\textwidth}
\centering
\includegraphics[width=\linewidth,height=4.5cm,keepaspectratio]{figures/terminale-11-1.pdf}
\end{column}
\end{columns}
\end{frame}
\begin{frame}{Activité de recherche}
\textbf{Encadrer une intégrale par des rectangles - 35 min.}

On cherche à encadrer \(I=\displaystyle\int_0^1x^2\,\dd x\). Pour un entier \(n\ge1\), on partage \([0;1]\) en \(n\) intervalles de même longueur. On note \(L_n\) la somme des aires des rectangles dont les hauteurs sont les valeurs de \(x^2\) aux extrémités gauches, et \(U_n\) celle obtenue aux extrémités droites.
\end{frame}
\begin{frame}{Activité de recherche (2)}
\begin{enumerate}
\setcounter{enumi}{0}\item Pour \(n=4\), représenter les deux familles de rectangles, puis calculer \(L_4\) et \(U_4\).

\item Justifier que \(L_4\le I\le U_4\).

\item Écrire une fonction Python d'argument \(n\) qui renvoie \(L_n\) et \(U_n\). Déterminer algébriquement \(U_n-L_n\).

\item Déterminer le plus petit entier \(n\) pour lequel cet encadrement a une largeur au plus égale à \(0{,}01\).

\item Calculer la valeur exacte de \(I\) à l'aide d'une primitive et vérifier qu'elle appartient à l'encadrement obtenu pour \(n=4\).

\end{enumerate}
\end{frame}
\end{document}
